\chapter{Regras de negócio e dados}\label{cap:regras}

\section{Catálogo de regras}
Uma regra expressa uma restrição do domínio. As fichas indicam os comportamentos
que a aplicam, e a matriz liga cada comportamento aos testes correspondentes.
As regras abaixo são hipóteses explícitas do produto fictício.

\begin{quadrosrs}
\caption{Regras e requisitos associados}
\begin{tabelasrs}{L{0.15\textwidth}YL{0.26\textwidth}}
\cabecalhosrs\textbf{Regra} & \textbf{Enunciado} & \textbf{Aplicação}\\\midrule
\phantomsection\label{rn:RN-001}\texttt{RN-001} & Uma pessoa tem no máximo uma inscrição ativa por evento; o total de inscrições ativas não supera a capacidade. & \reqref{RF-002}, \reqref{RF-003} e \reqref{RF-004}.\\
\phantomsection\label{rn:RN-002}\texttt{RN-002} & Criar ou cancelar uma inscrição exige instante estritamente anterior ao início do evento. A igualdade encerra o prazo. & \reqref{RF-002} e \reqref{RF-004}.\\
\end{tabelasrs}
\fonte{Elaboração didática deste guia.}
\end{quadrosrs}

A capacidade não pode ser reduzida para um valor inferior ao total de inscrições
ativas. Uma atualização externa que viole essa condição é rejeitada, preservando
os dados anteriores. A decisão de prazo e vaga usa o estado validado na operação
atômica de persistência.

\section{Dicionário de dados}
Identificadores são textos opacos, não vazios, com até 64 caracteres; não
transportam credenciais. A relação com participante usa somente a identidade
fornecida pelo serviço externo. Os nomes abaixo também orientam os contratos JSON.
Empates de início ou criação são desfeitos pela ordem lexicográfica dos
identificadores, comparando pontos de código Unicode, sem ignorar maiúsculas
ou aplicar regras de ordenação por idioma.

\begingroup\small\renewcommand{\arraystretch}{1.22}
\setlength{\tabcolsep}{5pt}
\begin{longtable}{L{0.29\textwidth}L{0.17\textwidth}L{0.12\textwidth}L{\dimexpr0.42\textwidth-8\tabcolsep\relax}}
\caption{Dados e restrições do domínio}\label{tab:dados}\\
\toprule\cabecalhosrs\textbf{Entidade e campo} & \textbf{Tipo lógico} & \textbf{Obrig.} & \textbf{Regra e significado}\\\midrule
\endfirsthead
\multicolumn{4}{l}{\textit{Dicionário de dados (continuação)}}\\
\toprule\cabecalhosrs\textbf{Entidade e campo} & \textbf{Tipo lógico} & \textbf{Obrig.} & \textbf{Regra e significado}\\\midrule
\endhead
\midrule\multicolumn{4}{r}{\textit{Continua na próxima página.}}\\\endfoot
\bottomrule\endlastfoot
Evento: \texttt{id} & Identificador & Sim & Único; recebido do catálogo externo.\\\hline
Evento: \texttt{nome} & Texto & Sim & Nome não vazio apresentado ao usuário.\\\hline
Evento: \texttt{inicio} & Instante UTC & Sim & Limite para criação e cancelamento.\\\hline
Evento: \texttt{capacidade} & Inteiro & Sim & Não negativo; não inferior às inscrições ativas.\\\hline
Evento: \texttt{publicado} & Booleano & Sim & Determina elegibilidade para catálogo e novas inscrições.\\\hline
Inscrição: \texttt{id} & Identificador & Sim & Único e imutável; gerado pelo produto.\\\hline
Inscrição: \texttt{eventoId} & Identificador & Sim & Referência a um evento conhecido.\\\hline
Inscrição: \texttt{participanteId} & Identificador & Sim & Titular proveniente da sessão; não informado pelo formulário.\\\hline
Inscrição: \texttt{situacao} & Enumeração & Sim & Valores \texttt{ativa} ou \texttt{cancelada}.\\\hline
Inscrição: \texttt{criadaEm} & Instante UTC & Sim & Momento da confirmação persistida.\\\hline
Inscrição: \texttt{canceladaEm} & Instante UTC & Cond. & Nulo quando ativa; preenchido no primeiro cancelamento.\\\hline
Evento: \texttt{vagas} & Inteiro derivado & Saída & Capacidade menos total de inscrições ativas; não negativo.\\
\end{longtable}
\endgroup

\section{Invariantes e consistência}
Cancelar conserva o registro e libera uma única vaga. Repetir a operação no
prazo preserva situação, instante do primeiro cancelamento e vagas. Uma nova
inscrição depois de cancelar recebe outro identificador e verifica novamente
capacidade, publicação e prazo.

Entradas inválidas e operações negadas não modificam registros nem contadores.
Falha de gravação não autoriza uma confirmação de sucesso. Se a resposta se
perder depois da confirmação persistida, a repetição da solicitação encontra
 a inscrição ativa e retorna seu identificador, conforme \reqref{RF-003}.

\section{Adequação ao domínio}
Na adaptação do modelo, identifique finalidade e necessidade de cada dado,
regras de retenção e exclusão, formatos de exportação, idiomas e obrigações
contratuais aplicáveis. Transforme obrigações do produto em requisitos com ID,
origem e aceite, e associe as regras às operações que as aplicam.
